\justifying
\NSFCBodyText

\subsubsection{与本项目相关的研究工作积累}

申请人长期致力于电力电子技术领域研究，\dots
近年来，申请人发表各类期刊论文\textbf{?篇}，\textbf{其中\textcolor{MsBlue}{第一作者}SCI期刊论文?篇（包含\textcolor{MsBlue}{中科院一区TOP期刊论文?篇）}}，\dots

申请人于2021年-2025年在\textbf{xx大学xx学院}攻读博士学位期间，师从\textcolor{MsBlue}{\textbf{xxx教授}}。
博士\dots

在本课题聚焦的\dots，申请人在以往的研究中已有很好的工作积累，取得了一定的研究成果。具体内容包括：

\textbf{\ding[1.2]{172}~~IPT系统\dots}~~
为解决\dots

\textbf{\ding[1.2]{173}~~适用于\dots}~~
为满足\dots

\subsubsection{已取得的研究工作成绩}

\ding[1.2]{172}~~\textbf{论文发表}：
在\textbf{IEEE TPEL\dots}等期刊发表\dots

\ding[1.2]{173}~~\textbf{专利申请}：
授权中国发明专利\dots

\ding[1.2]{174}~~\textbf{学科竞赛}：
分别于\dots

申请人近五年以\textbf{第一作者}发表的代表性论著如下：
\begin{enumerate}[label={[\arabic*]}, leftmargin=*, itemindent=0em, labelsep=0.5em]
	\item \textcolor{MsBlue}{\textbf{Ma Chunwei}}, Qu Xiaohui\textsuperscript{*}, Li Yundi, Liu Jinghang. Single-Stage Active Rectifier With Wide Impedance Conversion Ratio Range for Inductive Power Transfer System Delivering Constant Power. {\it \textbf{IEEE Transactions on Power Electronics}}, 2023, 38(6):7877 - 7890. （\textcolor{MsBlue}{\textbf{中科院一区TOP}}）
	\item \dots
\end{enumerate}

以上研究成果和所积累的理论\dots

\subsubsection{可行性分析}

\ding[1.2]{172}~~\textbf{理论可行性}：
本课题提出\dots

\ding[1.2]{173}~~\textbf{技术可行性}：本课题是申请人以往研究的延续。
\textbf{\dots 方面}：申请人前期\dots

\ding[1.2]{174}~~\textbf{实验条件可行性}：依托单位南通大学电气与自动化学院拥有先进的电力电子实验室\dots

\ding[1.2]{175}~~\textbf{前期验证}：申请人\dots

\subsubsection{研究风险的应对措施}

\subsubsubsection{动态\dots}

\textbf{风险描述}：在\dots。

\textbf{应对措施}：在\dots

\subsubsubsection{\dots 落地验证风险}

\textbf{风险描述}：本课题研究内容\dots

\textbf{应对措施}：针对\dots